\section{Major Examination, 2020--21 (Semester 2)}

\subsection*{Q1. True or false, with justification}

\begin{enumerate}
\item[(a)] \emph{While executing \texttt{char *c = (char *) malloc(30 *
sizeof(char));}, memory is allocated on the heap only.}

\textbf{False.}  The 30-byte object returned by \texttt{malloc} has allocated
storage duration and is normally obtained from the allocator's heap arena,
but the local pointer variable \texttt{c} has automatic storage duration
(normally a stack slot, unless optimized into a register or away).  Moreover,
a general-purpose allocator may obtain dynamic
storage using \texttt{mmap} rather than only by extending the traditional
\texttt{brk}-based heap.  Thus it is not correct to say that all memory
involved is allocated ``on the heap only.''

\item[(b)] \emph{A context switch can occur only after a timer interrupt.}

\textbf{False.}  A timer interrupt permits preemptive context switches, but
it is only one cause.  A running process may block in a system call, wait for
I/O, sleep, yield the processor, or incur a page fault that cannot be
resolved immediately.  In each case the scheduler may switch to another
runnable process.

\item[(c)] \emph{Threads \(T_1\) and \(T_2\) share \(x\).  \(T_1\) protects
\(x\) using a POSIX mutex.  If \(T_2\) writes \(x\) without locking the
mutex, a trap is generated.}

\textbf{False.}  A POSIX mutex is a software synchronization convention;
ordinary memory accesses do not carry information about which mutex should
protect them.  The hardware therefore does not trap merely because \(T_2\)
failed to lock the mutex.  The accesses constitute a data race (and hence
undefined behaviour in the C/POSIX memory model) if they are conflicting and
concurrent, but no protection trap is inherently generated.

\item[(d)] \emph{A thread cannot go to sleep while holding a lock, because
that might violate atomicity.}

\textbf{False.}  A thread can, for example, issue a blocking I/O operation
while retaining a mutex.  Other threads then remain excluded, so the mutual
exclusion, and hence the critical section's atomicity with respect to those
threads, is not violated.  Doing this is often poor design because it can
cause long delays, lock convoys, or deadlock.  Some special locks, such as
kernel spinlocks, impose a separate rule that their holders must not sleep,
but that is not a property of every lock.

\item[(e)] \emph{With a write-back disk buffer cache, a process writing a
block remains blocked until that block reaches the disk.}

\textbf{False.}  In a write-back cache the write normally modifies the
in-memory cached block and marks it dirty.  The call can return before a
later background write sends the block to disk.  An explicit operation such
as \texttt{fsync} is needed when the process requires the data to reach
stable storage before proceeding (and a write can still block for unrelated
resource-pressure reasons).
\end{enumerate}

\subsection*{Q2. Traps and limited direct execution}

\begin{enumerate}
\item[(a)] \emph{What are trap handlers?  Explain the hardware/software
interaction and the steps of trap handling.}

A trap handler is trusted kernel code that handles a transfer of control
caused by a system call, a synchronous exception (such as a protection fault
or divide-by-zero), or an asynchronous interrupt.  At boot, the kernel
installs handler addresses in a protected trap/interrupt vector table.

A typical trap proceeds as follows.
\begin{enumerate}
\item The process executes a trap-generating instruction, suffers an
exception, or is interrupted.
\item Hardware records the cause and saves enough state to resume execution,
such as the program counter and status register.  It changes from user mode
to kernel mode, selects the protected kernel stack, and vectors to the
registered entry point.  User code cannot choose an arbitrary kernel entry
address.
\item The low-level entry stub saves the remaining registers in a trap frame.
The kernel handler examines the cause, validates any system-call arguments,
and performs the requested service or handles the fault/interrupt.  It may
wake processes and invoke the scheduler.
\item To resume a process, the kernel restores its trap frame and executes a
privileged return-from-trap instruction.  This restores the saved program
counter and user privilege level (or restores a different process if a
context switch occurred).
\end{enumerate}

\item[(b)] \emph{How does trap handling help implement limited direct
execution (LDE)?}

LDE lets an application execute its ordinary instructions directly on the
CPU for speed, but only in user mode.  Whenever the application needs a
privileged service, it executes the designated system-call trap and enters
the kernel through a controlled entry point.  Faults transfer control to the
kernel when the application violates a protection rule, and interrupts let
the kernel regain control even if the application does not voluntarily trap.
The OS thus retains control while most application instructions run
natively.

\item[(c)] \emph{Give three forms of hardware support used by the OS to stop
a program from taking over the machine or doing something harmful in the
LDE model.}

\begin{enumerate}
\item \textbf{Privilege modes and privileged instructions:} user code cannot
disable interrupts, alter page tables, perform raw device I/O, or install
trap handlers.  A controlled trap instruction is used to request kernel
services.
\item \textbf{A programmable timer interrupt:} the OS arranges a future
interrupt before running a process, guaranteeing that it can regain the CPU
and preempt an infinite loop.
\item \textbf{Memory protection through the MMU:} page tables (or
base-and-bounds registers) restrict the process to authorized memory and
cause a protection fault on an illegal access.
\end{enumerate}
\end{enumerate}

\subsection*{Q3. \texttt{fork}, \texttt{exec}, and \texttt{wait}}

\begin{enumerate}
\item[(a)] \emph{Give all possible outputs of the following program, covering
both success and failure of system calls.  The command \texttt{wc} prints
the count for \texttt{p1.c}.}

\begin{verbatim}
int rc = fork();
if (rc == 0) {
    printf("Hello1\n");
    char *myargs[3];
    myargs[0] = strdup("wc");
    myargs[1] = strdup("p1.c");
    myargs[2] = NULL;
    execvp(myargs[0], myargs);
    printf("Hello2\n");
} else if (rc > 0) {
    wait();
    printf("Hello3\n");
} else {
    printf("Hello4\n");
}
\end{verbatim}

Treating the paper's \texttt{wait()} as \texttt{wait(NULL)}, and assuming
the displayed newline-terminated output is observed, the cases are:

\begin{enumerate}
\item If \texttt{fork} fails, there is no child and the output is
\begin{verbatim}
Hello4
\end{verbatim}

\item If \texttt{fork} and \texttt{execvp} both succeed, \texttt{execvp}
replaces the child, so \texttt{Hello2} is not executed.  The order is
\begin{verbatim}
Hello1
<wc output for p1.c>
Hello3
\end{verbatim}
The first line precedes \texttt{wc} by child program order, and
\texttt{Hello3} is last because the parent waits for the child.  If
\texttt{p1.c} does not exist, \texttt{execvp} has still succeeded:
\texttt{wc} prints an error diagnostic instead of a count, then terminates,
and \texttt{Hello3} follows.

\item If \texttt{fork} succeeds but \texttt{execvp} fails, execution
continues after \texttt{execvp} in the child:
\begin{verbatim}
Hello1
Hello2
Hello3
\end{verbatim}
Again, the wait ensures that the parent's line is last.
\end{enumerate}

\item[(b)] \emph{A parent forks a child that executes \texttt{cmd1}, then
forks a second child that executes \texttt{cmd2}, and finally calls
\texttt{wait} twice.  Is the order deterministic?  How can that behaviour
be reversed?}

The two commands execute \textbf{concurrently}, and their execution,
completion, and output order is non-deterministic.  The fact that the first
\texttt{fork} is called first does not force its child to run first.  Also,
the two waits occur only after both children have been created, so they
cannot impose an order between the commands.

To make the order deterministic, the parent can wait specifically for the
first child before creating the second:

\begin{verbatim}
pid_t p1 = fork();
if (p1 == 0) {
    exec(cmd1);
    _exit(1);                 /* reached only if exec fails */
}
waitpid(p1, NULL, 0);

pid_t p2 = fork();
if (p2 == 0) {
    exec(cmd2);
    _exit(1);
}
waitpid(p2, NULL, 0);
\end{verbatim}

This guarantees that \texttt{cmd1} terminates before \texttt{cmd2} starts.
Conversely, moving both waits after both forks, as in the question, changes
a sequential version back to non-deterministic concurrent execution.
\end{enumerate}

\subsection*{Q4. Six-level paging}

\emph{A byte-addressable system has 64-bit virtual addresses, 8-KiB pages,
and 8-byte PTEs.  It uses demand allocation at every page-table level.}

The page offset has
\[
 \log_2(8192)=13
\]
bits.  A page-table page contains
\[
 \frac{8192}{8}=1024=2^{10}
\]
PTEs, so a full level indexes 10 virtual-address bits.  The virtual page
number has \(64-13=51\) bits.

\begin{enumerate}
\item[(a)] \emph{How many page-table levels are required?}

\[
 L=\left\lceil\frac{51}{10}\right\rceil=6.
\]
Five levels use 10 bits each and the top level uses the remaining one bit.
This is a generic full 64-bit virtual address space; the 48-bit canonical
address restriction of some x86-64 implementations must not be imposed on
this question.

\item[(b)] \emph{Which virtual-address bits index each level?}

Numbering bits from 63 (most significant) to 0, one valid division is
\[
\underbrace{[63]}_{L_1}\;
\underbrace{[62:53]}_{L_2}\;
\underbrace{[52:43]}_{L_3}\;
\underbrace{[42:33]}_{L_4}\;
\underbrace{[32:23]}_{L_5}\;
\underbrace{[22:13]}_{L_6}\;
\underbrace{[12:0]}_{\mbox{page offset}} .
\]

\item[(c)] \emph{What is the maximum memory occupied by all pages of one
process's page table?}

For a completely populated 64-bit virtual address space, the numbers of
page-table pages from the root through the leaf level are
\[
1,\quad 2,\quad 2^{11},\quad 2^{21},\quad 2^{31},\quad 2^{41}.
\]
For example, there are \(2^{51}\) virtual data pages, and each leaf table
contains \(2^{10}\) PTEs, giving \(2^{41}\) leaf page-table pages.  Hence
\[
\begin{split}
N_{\rm PT}
 &= 1+2+2^{11}+2^{21}+2^{31}+2^{41}\\
 &=2{,}201{,}172{,}838{,}403\ \mbox{page-table pages}.
\end{split}
\]
The maximum memory is therefore
\[
N_{\rm PT}\,2^{13}
=18{,}032{,}007{,}892{,}197{,}376\ \mbox{bytes}
\mathrel{\simeq}16.0156\ \mbox{PiB}.
\]

\item[(d)] \emph{The process has accessed 8K (\(8192\)) data pages and none
has been swapped out.  What is the minimum memory needed for its page table?}

In the best case the virtual pages are packed into the fewest leaf tables.
Each leaf table maps 1024 data pages, so eight leaf tables suffice.  All
eight can be children of one \(L_5\) table and can share one table at each
of \(L_1,L_2,L_3,L_4\).  Thus
\[
N_{\min}=8+5=13\ \mbox{page-table pages},
\]
and
\[
M_{\min}=13(8\ \mbox{KiB})=104\ \mbox{KiB}
=106{,}496\ \mbox{bytes}.
\]

\item[(e)] \emph{For those 8192 accessed pages, what is the maximum memory
needed for all page-table entries in the worst placement?}

Page-table storage is allocated in whole 8-KiB pages, even if only one PTE
in such a page is valid.  Let \(P=8192\).  At each level, the number of
distinct table pages is bounded both by \(P\) and by the number of distinct
VPN prefixes that can reach that level:
\[
\begin{array}{c|c|c}
\mbox{level} & \mbox{maximum possible table pages} &
              \mbox{pages needed in the worst placement}\\ \hline
L_1 & 1       & 1\\
L_2 & 2       & 2\\
L_3 & 2^{11}  & 2048\\
L_4 & 2^{21}  & 8192\\
L_5 & 2^{31}  & 8192\\
L_6 & 2^{41}  & 8192
\end{array}
\]
These maxima can be attained simultaneously: use both top-level branches,
spread the pages over all 2048 available 11-bit prefixes, and choose a
different 21-bit (and hence 31-bit and 41-bit) prefix for every data page.
Consequently,
\[
\begin{split}
N_{\max}
 &=1+2+2048+3(8192)\\
 &=26{,}627\ \mbox{page-table pages},\\
M_{\max}
 &=26{,}627(8192)\\
 &=218{,}128{,}384\ \mbox{bytes}\\
 &=213{,}016\ \mbox{KiB}
 \mathrel{\simeq}208.0234\ \mbox{MiB}.
\end{split}
\]
In particular, \(8192\times6\) is not possible because there is only one
root table, at most two \(L_2\) tables, and at most 2048 \(L_3\) tables.
\end{enumerate}

\subsection*{Q5. Ticket lock}

\emph{Ten threads correctly use a ticket lock whose \texttt{lock} operation
atomically fetches and increments \texttt{ticket}, waits for
\(\texttt{turn}=\texttt{myturn}\), and whose \texttt{unlock} increments
\texttt{turn}.  Determine whether each given
\((\texttt{ticket},\texttt{turn})\) pair is possible.}

Ignoring integer wraparound, correct use maintains
\[
0\leq \texttt{ticket}-\texttt{turn}\leq 10.
\]
The difference is exactly the number of threads that have taken a ticket but
have not yet completed the corresponding unlock: one may be in the critical
section and the rest are waiting.

\begin{enumerate}
\item[(a)] \((0,0)\): \textbf{possible}.  This is the initialized state,
before any acquisition.

\item[(b)] \((0,1)\): \textbf{impossible}.  It would mean that an unlock had
advanced \texttt{turn} although no ticket had ever been issued.

\item[(c)] \((1,0)\): \textbf{possible}.  One thread has received ticket 0
and is either about to enter or is in the critical section; it has not yet
unlocked.

\item[(d)] \((150,100)\): \textbf{impossible}.  There would be
\(150-100=50\) outstanding acquisitions, but only ten correctly using
threads can each have at most one outstanding acquisition.

\item[(e)] \((1005,1000)\): \textbf{possible}.  One thousand acquisitions
have completed and five of the ten threads currently have outstanding
tickets.
\end{enumerate}

\subsection*{Q6. Racy linked-list insertion}

\emph{Three threads insert keys 1, 2, and 3 into an initially empty list,
without synchronization, using}

\begin{verbatim}
n->next = head;
head = n;
\end{verbatim}

\emph{where allocation always succeeds and is thread-safe.  Give every
possible reachable list length and shape.}

Under the intended interleaving model, pointer reads and writes are atomic.
Two or more threads can read the same old value of \texttt{head} in the
first assignment.  Their later writes to \texttt{head} then overwrite one
another, making some allocated nodes unreachable.

\begin{itemize}
\item Length 1:
\[
1,\qquad 2,\qquad 3.
\]
For example, all three threads can read \texttt{NULL} before any of them
writes \texttt{head}; whichever writes \texttt{head} last is the only
reachable node.

\item Length 2 (every ordered pair of distinct keys):
\[
\begin{array}{lll}
1\mathbin{\to}2,&1\mathbin{\to}3,&
2\mathbin{\to}1,\\
2\mathbin{\to}3,&3\mathbin{\to}1,&
3\mathbin{\to}2.
\end{array}
\]
The remaining node was allocated but lost.  For example, threads 2 and 3
may both prepare a node pointing to \texttt{NULL}; thread 2 publishes its
node, thread 1 records node 2 as its successor, thread 3 temporarily
overwrites the head, and thread 1 finally publishes node 1.  The reachable
list is \(1\to2\).

\item Length 3 (every permutation):
\[
\begin{array}{lll}
1\mathbin{\to}2\mathbin{\to}3,&
1\mathbin{\to}3\mathbin{\to}2,&
2\mathbin{\to}1\mathbin{\to}3,\\
2\mathbin{\to}3\mathbin{\to}1,&
3\mathbin{\to}1\mathbin{\to}2,&
3\mathbin{\to}2\mathbin{\to}1.
\end{array}
\]
Each is obtained by a suitable serial ordering (or an equivalent
interleaving) of the three insertions.
\end{itemize}

Length 0 is impossible because at least one of the three writes to
\texttt{head} occurs.  No cycle can be formed: an edge from node \(i\) to
node \(j\) requires \(j\) to have been published before \(i\)'s first
assignment; the resulting strict ``published earlier'' order cannot form a
cycle.  Strictly speaking, unsynchronized conflicting accesses to
\texttt{head} are a C data race and therefore have undefined behaviour;
the lists above are the complete outcomes in the sequentially consistent,
atomic-pointer model intended by the question.

\subsection*{Q7. Synchronizing a one-doctor clinic}

\emph{There is one doctor, an infinite waiting room, and one shared buffer.
A patient calls \texttt{consultDoctor()} to write consultation information
and later calls \texttt{noteTreatment()} to read the treatment.  The doctor
repeatedly calls \texttt{treatPatient()}, which reads the consultation and
writes the treatment.  Use only semaphores so patients do not overwrite one
another, the doctor reads only new valid information, and a patient reads
only a completed treatment.}

Three semaphores suffice:
\[
\begin{array}{c|c|l}
\mbox{name} & \mbox{initial value} & \mbox{purpose}\\ \hline
\texttt{office} & 1 & \mbox{only one patient may use the shared buffer}\\
\texttt{patientInfo} & 0 & \mbox{a complete, new consultation is available}\\
\texttt{treatmentReady} & 0 & \mbox{the doctor has written the treatment}
\end{array}
\]
Here \(P\) is semaphore wait and \(V\) is semaphore signal.

\begin{verbatim}
Patient()
{
    P(office);             /* wait outside; reserve the buffer */
    consultDoctor();       /* write this patient's information */
    V(patientInfo);        /* publish it to the doctor */

    P(treatmentReady);     /* wait until doctor writes result */
    noteTreatment();       /* read this patient's treatment */
    V(office);             /* next patient may use the buffer */
}

Doctor()
{
    while (1) {
        P(patientInfo);    /* never read an empty/stale request */
        treatPatient();    /* read request and write treatment */
        V(treatmentReady);
    }
}
\end{verbatim}

\texttt{office} is held by a patient for the complete request--response
exchange, so another patient cannot overwrite either phase of the shared
buffer.  The two zero-valued semaphores provide the request and response
handshakes, including the required memory-ordering points.  The semaphore's
wait queue represents the unbounded waiting room.

\subsection*{Q8. RAID and links}

\begin{enumerate}
\item[(a)] \emph{What is the steady-state random-write throughput of RAID-4,
and why?}

A small random write uses read--modify--write:
read the old data and old parity, compute
\[
P_{\rm new}=P_{\rm old}\mathbin{\oplus}D_{\rm old}
             \mathbin{\oplus}D_{\rm new},
\]
and then write the new data and new parity.  Thus every logical write makes
two accesses (one read and one write) to the single dedicated parity disk.
That disk is a serialization bottleneck, regardless of how many data disks
are present.

If one disk can sustain \(R\) random I/O operations per second, the precise
idealized small-random-write limit is about \(R/2\) logical writes per
second, because the parity disk performs two I/Os per logical write.  Thus
the throughput is of the order of one disk, not \(N-1\) disks.  (Caching or
full-stripe/coalesced writes can change this workload-specific figure.)

\item[(b)] \emph{How does RAID-5 improve random-write performance over
RAID-4?}

RAID-5 rotates parity blocks across all disks instead of placing all parity
on one dedicated disk.  A small write still has the four-I/O
read--modify--write penalty (old data read, old parity read, new data write,
new parity write), but unrelated writes can update parity blocks on
different disks.  With \(N\) disks each capable of \(R\) random IOPS, a
balanced idealized throughput is therefore approximately
\[
\frac{NR}{4}
\]
logical random writes per second, rather than RAID-4's parity-disk limit of
approximately \(R/2\).

\item[(c)] \emph{Distinguish hard links from soft (symbolic) links.}

\begin{center}
\begin{tabular}{p{0.43\linewidth}|p{0.43\linewidth}}
\textbf{Hard link} & \textbf{Symbolic (soft) link}\\ \hline
Another directory entry containing the same inode number as the original
name. & A distinct file whose contents are a pathname to another file.\\
All hard-linked names are equal names for the same inode; deleting one name
does not remove the file while another hard link remains. &
Deleting or moving the target can leave a dangling symbolic link.\\
Cannot cross file-system boundaries because inode numbers are local to a
file system. & Can contain a path that crosses file-system boundaries.\\
Hard links to directories are normally forbidden to preserve a well-formed
directory tree. & May refer to a directory.\\
Uses the target inode's metadata and increases its link count. &
Has its own inode; normal access follows the stored path to the target.
\end{tabular}
\end{center}
\end{enumerate}
